\section{Discussion and Limitations}
\label{sec:discussion}

\paragraph{What the evidence supports.}
Across two efficient-decoder seeds, the direction of the central observation is
stable: additional decoder capacity produces a small, bidirectional aggregate
effect, and its net benefit increases with an uncertainty proxy.  Boundary and
organ analyses further show that segmentation difficulty is spatially
heterogeneous.  Together, these findings justify testing selective regional
decoding rather than applying a uniformly expensive decoder.

\paragraph{What it does not yet support.}
The current study does not demonstrate lower \emph{executed} compute for an
adaptive model.  The legacy top-20\% result is a voxel oracle, not a realizable
router.  Even a block opportunity curve uses full-model predictions only to
define retrospective benefit; it does not account for errors introduced by a
learned refinement module.  Dynamic FLOPs, latency, memory, and final Dice must
therefore be evaluated in a separate method study.

\paragraph{Reproduction gap.}
Our efficient baseline remains 1.70--2.25 Dice points below the published
result, while the full baseline is much closer.  The most plausible unresolved
difference is data preprocessing: our reconstructed files have identity affine
metadata, whereas the original work used a processed dataset not distributed
with the repository.  The enlarged local full--efficient gap could inflate the
amount of apparent decoder opportunity.  We mitigate this by disclosing all
values, analyzing two efficient seeds, and limiting claims to controlled local
comparisons.  A paper-equivalent preprocessing run and matched multi-seed full
baseline remain important controls.

\paragraph{Statistical limitations.}
The dataset contains only 12 validation subjects.  Confidence intervals are
therefore wide, and millions of voxels must not be interpreted as independent
observations.  Our early entropy--error coefficient pooled bins across subjects
and is retained only as a descriptive diagnostic.  Final error predictability
will use subject-level AUROC/AUPRC and calibration.  Budgets and primary
endpoints must be declared before examining the corrected O9 results.

\paragraph{Generalization.}
All completed results currently use one CT dataset and one backbone family.
Without cross-backbone and cross-modality replication, the finding may be
specific to 3D UX-Net, the Synapse split, or EffiDec3D's resolution restriction.
The planned generality ladder separates matched decoder comparisons (3D UX-Net
and Swin UNETR) from ecological architecture controls (UNETR and nnU-Net), and
uses representative tasks from the original EffiDec3D scope (BTCV, FeTA, and
predeclared MSD lesion/thin-structure tasks).  The final claim will be narrowed
if these experiments do not reproduce the direction.

\paragraph{Clinical interpretation.}
High uncertainty and small-organ difficulty do not imply clinical safety.
Entropy can miss confidently wrong predictions, and improving average Dice
does not guarantee preservation of clinically important boundaries.  Any
future adaptive decoder should therefore retain a conservative fallback and be
evaluated per organ and per subject, not only by average compute.
